% BEGIN PROFESSOR FEEDBACK PF24: Discussion - Chapter introduction
\chapter{Discussion}
\label{ch:discussion}

This thesis examined whether temporal motifs provide useful information for
detecting malicious Ethereum addresses and for distinguishing between different
categories of malicious behaviour. The results lead to one overall conclusion.
Motif fingerprints contain information about address behaviour,
but they do not reliably improve classification once conventional transaction
and graph features are already available. The temporal window has a clear effect
on how many motifs can be observed, but the experiments do not identify one
unique window that is best for every sample. The Real-CATS categories also have
detectable differences in their fingerprints, although these differences are
too weak for accurate category classification. Selecting a separate window for
each category does not improve this result.

These findings are discussed below in the same order as the research questions.
The chapter then considers their theoretical and practical implications and the
limitations of the study.
% END PROFESSOR FEEDBACK PF24

